% ============================================================
%  数字逻辑设计课程设计实验报告
%  题目：基于 FPGA 的弹幕射击游戏（Danmaku Shooting Game）
%  平台：Sword Kintex-7 (XC7K160TFFG676-1)
%  开发环境：Xilinx Vivado 2025.1
% ============================================================

\documentclass[12pt,a4paper]{ctexart}

% ---------- 页面设置 ----------
\usepackage[top=2.5cm, bottom=2.5cm, left=2.8cm, right=2.5cm, headheight=14.5pt]{geometry}
\usepackage{graphicx}
\usepackage{float}
\usepackage{amsmath}
\usepackage{amssymb}
\usepackage{array}
\usepackage{booktabs}
\usepackage{longtable}
\usepackage{listings}
\usepackage{xcolor}
\usepackage{hyperref}
\usepackage{enumitem}
\usepackage{fancyhdr}
\usepackage{titlesec}
\usepackage{caption}
\usepackage{subcaption}
\usepackage{multirow}
\usepackage{makecell}
\usepackage{tikz}
\usetikzlibrary{automata, positioning, arrows.meta, shapes, chains, fit, backgrounds, shapes.geometric, shapes.multipart, calc, decorations.pathreplacing}

% ---------- 代码高亮 ----------
\lstset{
    basicstyle=\footnotesize\ttfamily,
    breaklines=true,
    frame=single,
    numbers=left,
    numberstyle=\tiny\color{gray},
    keywordstyle=\color{blue},
    commentstyle=\color{green!50!black},
    stringstyle=\color{red},
    tabsize=2,
    language=Verilog
}

% ---------- 页眉页脚 ----------
\pagestyle{fancy}
\fancyhf{}
\fancyhead[L]{数字逻辑设计课程设计实验报告}
\fancyhead[R]{浙江大学}
\fancyfoot[C]{\thepage}
\renewcommand{\headrulewidth}{0.4pt}
\renewcommand{\footrulewidth}{0pt}

% ---------- 标题格式 ----------
\titleformat{\section}{\Large\bfseries}{\thesection}{1em}{}
\titleformat{\subsection}{\large\bfseries}{\thesubsection}{1em}{}
\titleformat{\subsubsection}{\normalsize\bfseries}{\thesubsubsection}{1em}{}

\begin{document}

% ==================== 封面 ====================
\begin{titlepage}
\centering

\vspace*{3cm}

{\Huge\bfseries 数字逻辑设计\\课程设计实验报告\\}

\vspace{1.5cm}

{\LARGE 基于 FPGA 的弹幕射击游戏\\}

\vspace{0.8cm}

{\Large Danmaku Shooting Game on FPGA\\}

\vspace{2cm}

\begin{tabular}{rl}
\toprule
\textbf{课程名称} & 数字逻辑设计 \\
\textbf{实验平台} & Sword Kintex-7 (XC7K160TFFG676-1) \\
\textbf{开发环境} & Xilinx Vivado 2025.1 \\
\textbf{硬件描述语言} & Verilog HDL \\
\textbf{提交日期} & 2026 年 7 月 \\
\bottomrule
\end{tabular}

\vspace{2cm}

% TODO: 填写组内成员信息
\begin{tabular}{ccc}
\toprule
\textbf{学号} & \textbf{姓名} & \textbf{贡献比例} \\
\midrule
\multicolumn{1}{c}{\texttt{3210104669}} & {姜快快} & 100\% \\
\bottomrule
\end{tabular}

\vfill

\end{titlepage}

% ==================== 目录 ====================
\newpage
\tableofcontents
\newpage

% ==================== 一、背景介绍 ====================
\section{背景介绍}

本课程设计基于 Sword Kintex-7 实验平台（主控芯片 XC7K160TFFG676-1），
要求设计一个具有完整功能的时序电路，综合运用有限状态机、寄存器/存储器访问、
人机交互 I/O 接口等数字电路设计要素。

本小组选择设计一个基于 FPGA 的弹幕射击游戏。弹幕射击游戏（Danmaku Shooting）是一种经典的
视频游戏类型：玩家控制飞行器在屏幕上移动，发射子弹攻击敌方单位，同时躲避敌方子弹。
该选题兼具趣味性和技术挑战性，涵盖了 FSM 设计、VGA 显示控制、碰撞检测、
伪随机数生成、优先级合成渲染等数字电路设计核心内容。

本报告将详细介绍该游戏的设计思路、系统架构及各模块实现细节。

% ==================== 二、设计说明 ====================
\section{设计说明}

\subsection{设计开发环境}

\begin{itemize}
    \item \textbf{实验平台}：Sword Kintex-7 开发板（XC7K160TFFG676-1）
    \item \textbf{开发环境}：Xilinx Vivado 2025.1
    \item \textbf{硬件描述语言}：Verilog HDL
    \item \textbf{系统时钟}：100MHz（板载晶振，引脚 AC18）
    \item \textbf{显示输出}：VGA 640$\times$480 @ 60Hz，12 位色深（RGB 各 4 位）
    \item \textbf{输入设备}：PS/2 键盘（通过 Keyboard.v 模块解码）
    \item \textbf{辅助显示}：4 位 7 段数码管、8 位 LED
\end{itemize}

\subsection{输入输出交互选择}

\subsubsection{键盘输入}

\begin{itemize}
    \item \textbf{PS/2 键盘}：通过 \texttt{Keyboard.v} 模块解码 PS/2 Set 2 协议，输出 8 位按键状态 bitmap
    和 4 位难度选择 one-hot 信号。
    按键映射如下：
    \begin{itemize}
        \item W/S/A/D：上/下/左/右移动
        \item J：开火/开始游戏
        \item P：暂停/恢复
        \item 1/2/3/4：选择难度（EASY/NORMAL/HARD/HELL）
        \item M：切换模式（积分/无尽）
        \item Tab：切换 7 段数码管显示内容
    \end{itemize}
\end{itemize}

\subsubsection{显示输出}

\begin{itemize}
    \item \textbf{VGA 显示器}：640$\times$480 @ 60Hz，12 位 BGR 色深，实时渲染游戏画面
    \item \textbf{4 位 7 段数码管（直连模式）}：可切换显示击毁敌人数/玩家生命值/当前积分，
    通过 \texttt{DispNum} 模块动态扫描驱动
    \item \textbf{4 位 7 段数码管（串行 P2S 模式）}：通过 \texttt{Seg7Device} + \texttt{ShiftReg}
    串行移位输出级联扩展
    \item \textbf{8 位 LED}：积分高位溢出时用二进制表示高位部分；串行 P2S 扩展 LED
\end{itemize}

\subsection{系统总体架构}

系统自顶向下分为三层，采用"逻辑引擎 → 状态数据 → 渲染管线"的数据流架构：

\begin{enumerate}[label=(\arabic*)]
    \item \textbf{顶层互连层（game\_top）}：
    负责 32 位自由运行计数器（100MHz 递增）的时钟分频、PS/2 键盘解码接入、
    各子系统例化和 K7 引脚连接
    \item \textbf{核心功能层}：
    \begin{itemize}
        \item VGA 渲染管线（vga\_top）：
        由 25MHz VGA 时钟驱动 \texttt{vgac} 控制器产生扫描坐标 (row\_addr, col\_addr)；
        所有渲染器并行接收坐标，纯组合逻辑输出该像素的 hit 和颜色；
        优先级 MUX 合成最终像素；同时检测 vs 下降沿生成 60Hz frame\_tick 脉冲
        \item 游戏逻辑引擎（game\_logic）：
        运行于 60Hz frame\_tick 域，管理 5 状态 FSM、实体状态更新（玩家/敌机/子弹/障碍物）、
        逐对 AABB 碰撞检测、LFSR 伪随机数生成、武器升级、计分系统
    \end{itemize}
    \item \textbf{外设驱动层}：
     PS/2 键盘驱动（Keyboard）、7 段数码管直连显示（DispNum + SegmentDecoder + Seg7Decode + Seg7Remap）、
    7 段数码管串行 P2S（Seg7Device + ShiftReg）、LED 串行 P2S（ShiftReg）
\end{enumerate}

\begin{figure}[H]
\centering
\resizebox{\textwidth}{!}{%
\begin{tikzpicture}[
    node distance = 0.8cm,
    % 竖排文字：窄 text width 令中文逐字换行
    block/.style = {
        rectangle, draw = black, thick,
        text width = 1.1cm,
        align = center, font = \footnotesize, rounded corners = 2pt,
        inner sep = 3pt
    },
    wideblock/.style = {
        rectangle, draw = black, very thick,
        minimum width = 17.0cm, minimum height = 1.0cm,
        align = center, font = \small\bfseries, rounded corners = 3pt,
        fill = blue!8
    },
    ioblock/.style = {
        rectangle, draw = black, thick,
        text width = 1.1cm,
        align = center, font = \footnotesize, rounded corners = 2pt,
        fill = yellow!12, inner sep = 3pt
    },
    coreblock/.style = {
        rectangle, draw = black, thick,
        text width = 1.3cm,
        align = center, font = \footnotesize, rounded corners = 2pt,
        fill = red!10, inner sep = 3pt
    },
    vgablock/.style = {
        rectangle, draw = black, thick,
        text width = 1.3cm,
        align = center, font = \footnotesize, rounded corners = 2pt,
        fill = green!10, inner sep = 3pt
    },
    subblock/.style = {
        rectangle, draw = black, thin,
        text width = 1.1cm,
        align = center, font = \scriptsize, rounded corners = 2pt,
        inner sep = 3pt
    },
    arr/.style = {->, >=Stealth, thick},
    bidi/.style = {<->, >=Stealth, thick},
]

% ===== 第一层：输入 =====
\node[ioblock] (PS2) at (0, 0) {PS/2\\键\\盘};
\node[block, fill=orange!10] (KBD) at (3.0, 0)
    {Key-\\board\\PS/2\\Set2\\解码};
\draw[arr] (PS2) -- (KBD);

% ===== 第二层：game_top（宽横条）=====
\node[wideblock] (TOP) at (1.5, -3.0) {game\_top（FPGA 顶层模块）};
\draw[arr] (KBD.south) -- (KBD |- TOP.north);

% ===== 第三层：核心子系统（大幅拉开水平间距）=====
\node[coreblock] (LOGIC) at (-6.0, -6.0)
    {\textbf{game}\\\textbf{logic}\\游戏\\逻辑\\引擎};
\node[vgablock] (VTOP) at (3.0, -6.0)
    {\textbf{vga}\\\textbf{top}\\VGA\\渲染\\管线};
\node[block, fill=orange!10] (CLKDIV) at (10.0, -6.0)
    {clkdiv\\{}[31:0]\\时钟\\分频};

\draw[arr] (TOP.south) -- ++(0,-0.4) -| (LOGIC.north);
\draw[arr] (TOP.south) -- ++(0,-0.4) -| (VTOP.north);
\draw[arr] (TOP.south) -- ++(0,-0.4) -| (CLKDIV.north);

% ===== game_logic 子模块（x=-8 与 x=-4 两列）=====
\node[subblock, fill=red!6] (FSM)  at (-8.0, -10.0) {5状\\态\\FSM};
\node[subblock, fill=red!6] (ENT)  at (-4.0, -10.0) {实\\体\\管理};
\node[subblock, fill=red!6] (COL)  at (-8.0, -13.0) {碰\\撞\\检测};
\node[subblock, fill=red!6] (LFSR) at (-4.0, -13.0) {LFSR\\随\\机数};
\node[subblock, fill=red!6] (SCR)  at (-6.0, -16.0) {计分\\武器\\升级};

\draw[arr, shorten >=2pt] (LOGIC.south) -- ++(0,-0.3) -| (FSM.north);
\draw[arr, shorten >=2pt] (LOGIC.south) -- ++(0,-0.3) -| (ENT.north);
\draw[arr] (FSM) -- (COL);
\draw[arr] (ENT) -- (LFSR);
\draw[arr] (COL.south) -- ++(0,-0.3) -| (SCR.north);

% ===== vga_top 子模块（x=1 与 x=5 两列）=====
\node[subblock, fill=green!6] (VGAC)  at (1.0, -10.0) {vgac\\控\\制器};
\node[subblock, fill=green!6] (REND)  at (5.0, -10.0) {渲\\染器\\阵列};
\node[subblock, fill=green!6] (MUX)   at (1.0, -13.0) {优先\\级\\MUX};
\node[subblock, fill=green!6] (FTICK) at (5.0, -13.0) {frame\\tick\\生成};

\draw[arr, shorten >=2pt] (VTOP.south) -- ++(0,-0.3) -| (VGAC.north);
\draw[arr, shorten >=2pt] (VTOP.south) -- ++(0,-0.3) -| (REND.north);
\draw[arr] (VGAC) -- (MUX);
\draw[arr] (REND) -- (FTICK);

% ===== 双向数据总线 =====
\draw[bidi, thick, color=red!70!black]
    (LOGIC.east) -- node[above, font=\footnotesize, text=red!70!black] {实体状态总线} (VTOP.west);

% ===== 第四层：输出设备 =====
\node[ioblock] (VGA)  at (1.0, -17.0) {VGA\\显示\\器\\640\\×480};
\node[ioblock] (SEG7) at (-2.0, -17.0) {7段\\数\\码管};
\node[ioblock] (LED)  at (-4.5, -17.0) {LED\\指\\示灯};
\node[ioblock] (P2S)  at (5.0, -17.0) {串行\\P2S\\输出};

\draw[arr, dashed] (MUX.south) -- ++(0,-0.4) -| (VGA.north);
\draw[arr, dashed] (SCR.south) -- ++(0,-0.4) -| (SEG7.north);
\draw[arr, dashed] (SCR.south) -- ++(0,-0.4) -| (LED.north);
\draw[arr, dashed] (FTICK.south) -- ++(0,-0.4) -| (P2S.north);

\end{tikzpicture}%
}
\caption{系统总体架构框图}
\label{fig:arch}
\end{figure}

\subsection{引脚约束}

系统通过 \texttt{K7.xdc} 文件定义所有 FPGA 引脚约束，包括：
\begin{itemize}
    \item \textbf{系统时钟}：100MHz 差分时钟，引脚 AC18（clk）
    \item \textbf{全局复位}：低有效复位，引脚 W13（rstn）
    \item \textbf{VGA 输出}：
    R[3:0]（M17/N16/L16/K17）、G[3:0]（J19/J20/K19/K20）、B[3:0]（M20/M19/L17/L18）、HS（M22）、VS（M21）
    \item \textbf{键盘输入}：PS/2 接口（ps2\_clk / ps2\_data）
    \item \textbf{7 段数码管}：SEGMENT[7:0]、AN[3:0]
    \item \textbf{串行输出}：SEG\_CLK/SEG\_DT/SEG\_EN/SEG\_CLR、LED\_CLK/LED\_DT/LED\_EN/LED\_CLR
    \item \textbf{LED}：LED[7:0]
\end{itemize}

\subsection{时钟域设计}

系统在 100MHz 主时钟下运行。关键时钟信号：

\begin{itemize}
    \item \textbf{vga\_clk = 25MHz}：2 位计数器 \texttt{cdiv[1:0]} 对 100MHz 进行 4 分频产生，
    驱动 VGA 控制器和所有纯组合渲染器
    \item \textbf{frame\_tick（约 60Hz）}：在 100MHz 域中通过三级同步器（\texttt{vs\_sync[2:0]}）
    检测 vs 信号的下降沿（\texttt{vs\_sync[2] \&\& !vs\_sync[1]}），每个 V-blank 期间产生一个周期的脉冲，
    驱动游戏逻辑状态更新
    \item \textbf{PS/2 键盘接口}：标准 PS/2 Set 2 协议解码
    \item \textbf{数码管动态扫描（约 6.25MHz）}：\texttt{clkdiv[15:14]} 用于 DispNum 的扫描选择信号
    \item \textbf{串行移位时钟}：\texttt{clkdiv[3]}（100MHz$\div 16 \approx 6.25$MHz）
    用于 ShiftReg 的串行输出
\end{itemize}


% ==================== 三、模块结构 ====================
\section{模块结构}

\subsection{工程文件}

工程源码共包含 18 个 Verilog 源文件（.v）和 1 个头文件（.vh），按功能分类如下：

\begin{longtable}{llp{5.5cm}}
\caption{工程源文件列表}
\label{tab:files} \\
\toprule
\textbf{类别} & \textbf{文件名} & \textbf{功能} \\
\midrule
\endfirsthead
\multicolumn{3}{c}{\tablename\ \thetable{}（续）} \\
\toprule
\textbf{类别} & \textbf{文件名} & \textbf{功能} \\
\midrule
\endhead
\midrule
\multicolumn{3}{r}{\footnotesize 接下页} \\
\endfoot
\bottomrule
\endlastfoot
头文件 & \texttt{game\_defs.vh} & 全局宏定义（屏幕尺寸、颜色、状态编码、实体数量、游戏参数等） \\
\midrule
\multirow{2}{*}{顶层} & \texttt{game\_top.v} & FPGA 顶层（时钟分频、PS/2 键盘接入、子系统互联、引脚映射） \\
                       & \texttt{DispNum.v} & 4 位 7 段数码管动态扫描显示驱动（含 2-4 译码器、多路选择器、MC14495 兼容译码器） \\
\midrule
VGA 基础 & \texttt{vgac.v} & ZJU VGA 控制器（640$\times$480@60Hz，输出扫描坐标、同步信号、消隐） \\
          & \texttt{vga\_top.v} & VGA 渲染管线顶层（25MHz 时钟分频、渲染器例化、优先级 MUX、frame\_tick 生成） \\
\midrule
\multirow{5}{*}{渲染器} & \texttt{player\_render.v} & 玩家三角形飞机渲染（纯组合逻辑） \\
                         & \texttt{enemy\_render.v} & 敌机椭圆 UFO 渲染（8 实例，generate 展开） \\
                         & \texttt{bullet\_render.v} & 子弹渲染（16 玩家菱形 + 64 敌弹方块，generate 展开） \\
                         & \texttt{menu\_render.v} & 菜单界面渲染（标题/装饰边框/难度/模式/操作提示文字） \\
                         & \texttt{hud\_render.v} & HUD 渲染（分数/生命/暂停/GAMEOVER/YOU WIN 文字） \\
\midrule
游戏逻辑 & \texttt{game\_logic.v} & 核心逻辑（5 状态 FSM、实体管理、碰撞检测、LFSR、计分、武器升级） \\
\midrule
\multirow{5}{*}{外设驱动} & \texttt{Keyboard.v} & PS/2 键盘驱动（Set 2 协议 make/break 解码 + 扩展键码 E0 前缀处理） \\
                        & \texttt{SegmentDecoder.v} & 单 digit 十六进制→7 段译码（低有效） \\
                        & \texttt{Seg7Decode.v} & 8 位 7 段译码 + 锁存使能（级联 8 个 SegmentDecoder） \\
                        & \texttt{Seg7Remap.v} & 7 段引脚重映射（补偿 PCB 布线交叉） \\
                        & \texttt{Seg7Device.v} & 7 段设备串行 P2S 输出（Seg7Decode + Seg7Remap + ShiftReg + 动态扫描） \\
                        & \texttt{ShiftReg.v} & 参数化并转串移位寄存器（LED P2S 输出） \\
\end{longtable}

注：\texttt{Keyboard.v} 已实例化到 \texttt{game\_top} 顶层中，
所有输入通过 PS/2 键盘完成。

\subsection{各模块关系图}

\begin{figure}[H]
\centering
\resizebox{\textwidth}{!}{%
\begin{tikzpicture}[
    node distance = 0.7cm,
    % 竖排文字：窄 text width 令中文逐字换行
    root/.style = {
        rectangle, draw = black, very thick,
        text width = 1.6cm,
        align = center, font = \small\bfseries, rounded corners = 3pt,
        fill = blue!10, inner sep = 4pt
    },
    branch/.style = {
        rectangle, draw = black, thick,
        text width = 1.2cm,
        align = center, font = \footnotesize, rounded corners = 2pt,
        fill = orange!10, inner sep = 3pt
    },
    leaf/.style = {
        rectangle, draw = black, thin,
        text width = 1.1cm,
        align = center, font = \footnotesize, rounded corners = 2pt,
        fill = gray!5, inner sep = 3pt
    },
    subleaf/.style = {
        rectangle, draw = black, thin,
        text width = 1.1cm,
        align = center, font = \scriptsize, rounded corners = 2pt,
        fill = gray!3, inner sep = 3pt
    },
    vleaf/.style = {
        rectangle, draw = black, thin,
        text width = 1.1cm,
        align = center, font = \footnotesize, rounded corners = 2pt,
        fill = green!8, inner sep = 3pt
    },
    subvleaf/.style = {
        rectangle, draw = black, thin,
        text width = 1.1cm,
        align = center, font = \scriptsize, rounded corners = 2pt,
        fill = green!4, inner sep = 3pt
    },
    gleaf/.style = {
        rectangle, draw = black, thin,
        text width = 1.1cm,
        align = center, font = \footnotesize, rounded corners = 2pt,
        fill = red!8, inner sep = 3pt
    },
    subgleaf/.style = {
        rectangle, draw = black, thin,
        text width = 1.1cm,
        align = center, font = \scriptsize, rounded corners = 2pt,
        fill = red!4, inner sep = 3pt
    },
    arr/.style = {-, draw=black, thick},
]

\node[root] (TOP) at (0, 0) {game\\top\\FPGA\\顶层};

% 两大分支
\node[branch] (CLK)  at (-5.0, -3.5) {clkdiv\\{}[31:0]\\时钟\\分频};
\node[branch] (KBD)  at (6.0, -3.5)  {Key-\\board\\PS/2\\驱动};

% game_logic / vga_top（逻辑与渲染）
\node[gleaf] (LOGIC) at (-9.0, -7.0)
    {\textbf{game}\\\textbf{logic}\\核心\\逻辑};
\node[vleaf] (VT)    at (-2.0, -7.0)
    {\textbf{vga}\\\textbf{top}\\渲染\\管线};

% DispNum / Seg7Device（外设）
\node[leaf] (DISP)   at (3.0, -7.0)  {Disp-\\Num\\直连\\7段};
\node[leaf] (SEGDEV) at (9.0, -7.0)  {Seg7\\Device\\串行\\P2S};

% game_logic 子模块
\node[subgleaf] (FSM)   at (-11.0, -11.0) {5状\\态\\FSM};
\node[subgleaf] (ENTS)  at (-8.0, -11.0)  {实\\体\\管理};
\node[subgleaf] (COLL)  at (-11.0, -14.0) {AABB\\碰撞\\检测};
\node[subgleaf] (RNG)   at (-8.0, -14.0)  {Galois\\LFSR};
\node[subgleaf] (SCORE) at (-9.5, -17.0)  {计分\\武器\\升级};

% vga_top 子模块
\node[subvleaf] (VG)      at (-4.0, -11.0) {vgac\\640\\×480};
\node[subvleaf] (RENDERS) at (-0.5, -11.0) {渲染\\器\\阵列};
\node[subvleaf] (PRIO)    at (-4.0, -14.0) {优先\\级\\MUX};
\node[subvleaf] (FT)      at (-0.5, -14.0) {frame\\tick\\同步};
\node[subvleaf, font=\tiny] (RPL) at (-0.5, -16.6) {player\\render\\×1};
\node[subvleaf, font=\tiny] (REN) at (-0.5, -18.6) {enemy\\render\\×8};
\node[subvleaf, font=\tiny] (RPB) at (-0.5, -20.6) {bullet\\render\\×80};
\node[subvleaf, font=\tiny] (RHU) at (-0.5, -22.6) {hud+\\menu\\render};

% 外设子模块
\node[subleaf] (SEGDEC)  at (3.0, -11.0) {Seg-\\ment\\Dec\\×8};
\node[subleaf] (SEGRMP)  at (9.0, -11.0) {Seg7\\Decode\\+Remap};
\node[subleaf] (SHIFT)   at (9.0, -14.0) {Shift\\Reg\\并转\\串};

\draw[arr] (TOP) -| (CLK);
\draw[arr] (TOP) -| (KBD);
\draw[arr] (CLK.south) -- ++(0,-0.4) -| (LOGIC.north);
\draw[arr] (CLK.south) -- ++(0,-0.4) -| (VT.north);
\draw[arr] (KBD.south) -- ++(0,-0.4) -| (DISP.north);
\draw[arr] (KBD.south) -- ++(0,-0.4) -| (SEGDEV.north);

\draw[arr] (LOGIC.south) -- ++(0,-0.3) -| (FSM.north);
\draw[arr] (LOGIC.south) -- ++(0,-0.3) -| (ENTS.north);
\draw[arr] (FSM) -- (COLL);
\draw[arr] (ENTS) -- (RNG);
\draw[arr] (COLL.south) -- ++(0,-0.3) -| (SCORE.north);

\draw[arr] (VT.south) -- ++(0,-0.3) -| (VG.north);
\draw[arr] (VT.south) -- ++(0,-0.3) -| (RENDERS.north);
\draw[arr] (VG) -- (PRIO);
\draw[arr] (RENDERS) -- (FT);
\draw[arr] (FT) -- (RPL);
\draw[arr] (RPL) -- (REN);
\draw[arr] (REN) -- (RPB);
\draw[arr] (RPB) -- (RHU);

\draw[arr] (DISP) -- (SEGDEC);
\draw[arr] (SEGDEV) -- (SEGRMP);
\draw[arr] (SEGRMP) -- (SHIFT);

\draw[<->, >=Stealth, draw=red!70!black, thick]
    (LOGIC.east) -- node[midway, above, font=\footnotesize, text=red!70!black] {实体状态总线} (VT.west);

\begin{scope}[on background layer]
\node[draw=red!50, dashed, thick, inner sep=8pt,
      fit=(LOGIC)(FSM)(COLL)(SCORE)(ENTS)(RNG),
      label=above:{\scriptsize\color{red!60} 逻辑引擎层}] {};
\node[draw=green!50!black, dashed, thick, inner sep=8pt,
      fit=(VT)(VG)(RENDERS)(PRIO)(FT)(RPL)(REN)(RPB)(RHU),
      label=above:{\scriptsize\color{green!50!black} 渲染管线层}] {};
\node[draw=orange!60!black, dashed, thick, inner sep=8pt,
      fit=(KBD)(DISP)(SEGDEV)(SEGDEC)(SEGRMP)(SHIFT),
      label=above:{\scriptsize\color{orange!60!black} 外设驱动层}] {};
\end{scope}

\end{tikzpicture}%
}
\caption{模块关系结构图}
\label{fig:modules}
\end{figure}

\subsection{核心模块详细说明}

\subsubsection{game\_top --- FPGA 顶层模块}

\texttt{game\_top} 是 FPGA 顶层模块，端口与 \texttt{K7.xdc} 引脚约束一一对应：

\begin{lstlisting}
module game_top (
    input  wire        clk,         // 100MHz 系统时钟
    input  wire        rstn,        // 低有效复位
    output wire [3:0]  r, g, b,     // VGA 色彩通道
    output wire        hs, vs,      // VGA 行/场同步
    input  wire        ps2_clk,     // PS/2 键盘时钟
    input  wire        ps2_data,    // PS/2 键盘数据
    output wire [7:0]  SEGMENT,     // 7 段数码管段选（直连）
    output wire [3:0]  AN,          // 7 段数码管位选（直连）
    output wire        SEG_CLK, SEG_CLR, SEG_DT, SEG_EN,  // 7 段串行 P2S
    output wire [7:0]  LED,         // LED 直连
    output wire        LED_CLK, LED_CLR, LED_DT, LED_EN   // LED 串行 P2S
);
\end{lstlisting}

顶层内部结构：
\begin{itemize}
    \item 32 位自由运行计数器 \texttt{clkdiv}（100MHz 递增），用于多级分频
    \item 例化 \texttt{Keyboard}（PS/2 键盘驱动，输出按键 bitmap、难度、模式等状态信号）
    \item 例化 \texttt{vga\_top}（VGA 渲染管线，输出 r/g/b/hs/vs 和 frame\_tick）
    \item 例化 \texttt{game\_logic}（核心游戏逻辑，接收 frame\_tick 和键盘状态，输出所有实体状态）
    \item 例化 \texttt{DispNum}（4 位 7 段数码管直连动态扫描显示）
    \item 例化 \texttt{Seg7Device} + \texttt{ShiftReg}（串行 7 段 P2S 和 LED P2S 输出）
\end{itemize}

\begin{figure}[H]
\centering
% TODO: 在此插入 game_top 模块的 Vivado RTL Schematic 截图
% 操作步骤：Vivado → Open Elaborated Design → RTL Analysis → Schematic → 选中 game_top → Show Hierarchy
\fbox{\parbox{0.85\textwidth}{\centering\vspace{3cm}【图：game\_top 模块 RTL Schematic】\vspace{3cm}}}
\caption{game\_top 模块 RTL Schematic}
\label{fig:gametop_rtl}
\end{figure}

\subsubsection{vgac --- VGA 控制器}

\texttt{vgac} 为标准 640$\times$480@60Hz VGA 时序控制器，驱动 ZJU VGA 接口。关键时序参数：

\begin{itemize}
    \item 水平时序：一行共计 800 像素时钟周期（h\_count 0$\sim$799），有效显示区 144$\sim$783（640 像素），行同步 h\_count $>$ 95 时输出高
    \item 垂直时序：一帧共计 525 行（v\_count 0$\sim$524），有效显示区 35$\sim$514（480 行），场同步 v\_count $>$ 1 时输出高
    \item 地址生成：row\_addr = v\_count $-$ 35（范围 0$\sim$479），col\_addr = h\_count $-$ 143（范围 0$\sim$639）
    \item 消隐控制：有效显示区内 rdn = 0（读使能），消隐期间 rdn = 1 强制输出黑色
    \item 颜色位序：\texttt{d\_in = \{B[3:0], G[3:0], R[3:0]\}}（BGR 序，非 RGB 序）
\end{itemize}

\subsubsection{vga\_top --- VGA 渲染管线顶层}

\texttt{vga\_top} 负责将游戏数据转换为 VGA 视频信号，是整个渲染系统的顶层：

\textbf{VGA 时钟}：通过 2 位计数器 \texttt{cdiv[1:0]} 对 100MHz 系统时钟 4 分频得到 25MHz 像素时钟。

\textbf{帧同步脉冲生成}：通过三级同步器 \texttt{vs\_sync[2:0]} 在 100MHz 域中检测 vs 信号的下降沿（1→0），
在 V-blank 期间产生一个周期的 \texttt{frame\_tick} 脉冲，频率约 60Hz。这一设计确保了游戏状态更新不会造成画面撕裂。

\textbf{渲染器阵列}：所有渲染器均为纯组合逻辑，在同一 25MHz 像素时钟域中并行工作。
每个渲染器接收扫描坐标 (sx, sy)，输出 hit（当前像素是否属于该实体）和 12 位 BGR 颜色信号。
渲染器实例通过 \texttt{generate} 块展开：

\begin{table}[H]
\centering
\caption{渲染器实例列表}
\label{tab:renderers}
\begin{tabular}{lccp{4cm}}
\toprule
\textbf{渲染器} & \textbf{实例数} & \textbf{展开方式} & \textbf{说明} \\
\midrule
player\_render & 1 & 单实例 & 玩家三角形飞机 \\
enemy\_render & 8 & generate for & 敌人椭圆 UFO，每个独立坐标和 HP \\
bullet\_render（玩家） & 16 & generate for & 玩家菱形子弹，is\_player=1 \\
bullet\_render（敌方） & 64 & generate for & 敌方红色方块子弹，is\_player=0 \\
hud\_render & 1 & 单实例 & HUD 文字叠加层（分数/HP/暂停/结束） \\
menu\_render & 1 & 单实例 & 菜单界面（仅在 STATE\_MENU 下有效） \\
\bottomrule
\end{tabular}
\end{table}

\textbf{优先级多路选择器（Priority MUX）}：通过级联的 \texttt{if-else} 优先级链实现像素合成。
消隐期间（rdn=1）强制输出黑色。其余层级按以下顺序从高到低：

\begin{center}
\small
Menu（仅 MENU 状态）＞ HUD ＞ 玩家飞机 ＞ 玩家子弹 ＞ 敌方子弹 ＞ 敌人 ＞ 黑色背景
\end{center}

优先级 MUX 核心代码：
\begin{lstlisting}
always @(*) begin
    if (rdn)                        vga_data = `COL_BLACK;
    else if (hit_menu && game_state == `STATE_MENU) vga_data = col_menu;
    else if (hit_hud)               vga_data = col_hud;
    else if (hit_pl)                vga_data = col_pl;
    else if (hit_pb)                vga_data = col_pb;
    else if (hit_eb)                vga_data = col_eb;
    else if (hit_en)                vga_data = col_en;
    else if (hit_ob)                vga_data = col_ob;
    else                            vga_data = `COL_BLACK;
end
\end{lstlisting}

\begin{figure}[H]
\centering
% TODO: 在此插入 vga_top 模块的 Vivado RTL Schematic 截图
% 操作步骤：Vivado → Open Elaborated Design → RTL Analysis → Schematic → 选中 vga_top
\fbox{\parbox{0.85\textwidth}{\centering\vspace{3cm}【图：vga\_top 模块 RTL Schematic】\vspace{3cm}}}
\caption{vga\_top 模块 RTL Schematic}
\label{fig:vgatop_rtl}
\end{figure}

\subsubsection{game\_logic --- 核心游戏逻辑}

\texttt{game\_logic} 是整个游戏的核心模块，由 frame\_tick（约 60Hz）驱动，在 100MHz 时钟域中运行。

\paragraph{游戏状态机}

系统采用 5 状态 Moore 型有限状态机，状态编码定义于 \texttt{game\_defs.vh}：

\begin{itemize}
    \item \textbf{STATE\_MENU}（3'd0）：主菜单界面，允许选择难度和游戏模式，等待开始信号
    \item \textbf{STATE\_PLAYING}（3'd1）：游戏进行中，实体状态更新、碰撞检测、计分
    \item \textbf{STATE\_PAUSED}（3'd2）：暂停状态，画面冻结，约 4 秒冷却后允许恢复
    \item \textbf{STATE\_GAMEOVER}（3'd3）：游戏结束（玩家生命值归零），显示 300 帧（5 秒）后返回 MENU
    \item \textbf{STATE\_WIN}（3'd4）：通关（积分模式达到 500 分），显示 300 帧（5 秒）后返回 MENU
\end{itemize}

\textbf{状态转换条件：}

\begin{enumerate}[label=(\arabic*)]
    \item MENU $\rightarrow$ PLAYING：在 MENU 状态下检测到 start\_key（J/Space/Enter 键）有效，
    记录当前选择的难度和模式，初始化玩家 HP（默认为 3，HELL 难度为 2），
    设置 \texttt{state\_timer = GRACE\_PERIOD（300 帧，约 5 秒）}开局无敌保护期
    \item PLAYING $\rightarrow$ PAUSED：在 PLAYING 状态下，检测到 pause\_key（P 键上升沿）
    且 \texttt{pause\_cooldown == 0} 时进入暂停。进入后设置 \texttt{pause\_cooldown = 250}（约 4 秒冷却）
    \item PAUSED $\rightarrow$ PLAYING：在 PAUSED 状态下，检测到 pause\_key 或 fire\_key 且
    \texttt{pause\_cooldown == 0} 时恢复游戏，同时设置冷却防止连按
    \item PLAYING $\rightarrow$ GAMEOVER：在 PLAYING 状态下，\texttt{pl\_hp == 0}（玩家生命归零）时触发。
    设置 \texttt{state\_timer = GAMEOVER\_DISPLAY（300 帧 = 5 秒）}
    \item PLAYING $\rightarrow$ WIN：在 PLAYING 状态下，积分模式下 \texttt{total\_score >= SCORE\_WIN\_TARGET（500）}
    且 \texttt{pl\_alive} 为真时触发。设置 \texttt{state\_timer = WIN\_DISPLAY（300 帧 = 5 秒）}
    \item GAMEOVER $\rightarrow$ MENU：\texttt{state\_timer} 倒计时至 0 后自动返回，同时清除所有实体
    \item WIN $\rightarrow$ MENU：\texttt{state\_timer} 倒计时至 0 后自动返回，同时清除所有实体
\end{enumerate}

\begin{figure}[H]
\centering
\resizebox{\textwidth}{!}{%
\begin{tikzpicture}[
    ->,
    >=Stealth,
    node distance = 5.0cm and 10.0cm,
    every state/.style = {
        draw = black,
        thick,
        fill = gray!10,
        minimum size = 2.0cm,
        font = \small,
        align = center
    },
    every edge/.style = {
        draw = black,
        thick
    }
]

% 五个状态节点（水平间距大幅拉开，纵向拉长）
\node[state] (MENU)     at (0, 0)       {MENU\\[-2pt] \scriptsize 3'd0};
\node[state] (PLAYING)  at (0, -7.0)    {PLAY-\\ING\\[-2pt] \scriptsize 3'd1};
\node[state] (PAUSED)   at (11.0, -7.0) {PAU-\\SED\\[-2pt] \scriptsize 3'd2};
\node[state] (GAMEOVER) at (-9.0, -14.0){GAME\\OVER\\[-2pt] \scriptsize 3'd3};
\node[state] (WIN)      at (9.0, -14.0) {WIN\\[-2pt] \scriptsize 3'd4};

% 1. MENU → PLAYING
\path (MENU)  edge node[right, font=\footnotesize] {start\_key (J)}  (PLAYING);

% 2. PLAYING → PAUSED
\path (PLAYING) edge [bend left = 12]
      node[above, font=\footnotesize] {P 键 \& cooldown=0} (PAUSED);

% 3. PAUSED → PLAYING
\path (PAUSED) edge [bend left = 12]
      node[below, font=\footnotesize] {P $\|$ J \& cooldown=0} (PLAYING);

% 4. PLAYING → GAMEOVER
\path (PLAYING) edge [bend right = 18]
      node[above left, font=\footnotesize, xshift=-0.3cm] {pl\_hp\,$=$\,0} (GAMEOVER);

% 5. PLAYING → WIN
\path (PLAYING) edge [bend left = 18]
      node[above right, font=\footnotesize, xshift=0.3cm] {score $\ge$ 500 \& !mode} (WIN);

% 6. GAMEOVER → MENU
\path (GAMEOVER) edge [bend right = 35]
      node[left, font=\footnotesize, align=center, xshift=-0.3cm] {timer=0\\(300 帧后)} (MENU);

% 7. WIN → MENU
\path (WIN) edge [bend left = 35]
      node[right, font=\footnotesize, align=center, xshift=0.3cm] {timer=0\\(300 帧后)} (MENU);

% 初始状态标记
\node[above = 0.3cm of MENU, font = \small\bfseries] {rstn 复位 $\to$ MENU};

\end{tikzpicture}%
}
\caption{游戏状态机状态转换图（Moore 型 FSM）}
\label{fig:fsm}
\end{figure}

\paragraph{输入映射}

\texttt{game\_logic} 通过 \texttt{Keyboard.v} 模块获取解码后的键盘输入信号：

\begin{table}[H]
\centering
\caption{键盘按键映射表}
\label{tab:input}
\begin{tabular}{lll}
\toprule
\textbf{按键} & \textbf{功能} & \textbf{类型} \\
\midrule
W / S / A / D  & 上 / 下 / 左 / 右移动    & 持续有效 \\
J        & 开火 / 开始游戏   & 持续有效 \\
P        & 暂停 / 恢复        & 上升沿锁存 \\
1/2/3/4  & 选择难度（EASY/NORMAL/HARD/HELL） & 即时生效 \\
M        & 切换模式（积分/无尽）  & 上升沿锁存 \\
Tab      & 切换 7 段显示内容    & 上升沿锁存 \\
\bottomrule
\end{tabular}
\end{table}

难度和模式选择在 MENU 状态下通过按键 1-4 和 M 键设定。武器升级为自动触发，
无需按 K 键：当击杀数达到 10/20/40... 时自动升级武器等级（最多 3 级）。

\paragraph{实体管理}

采用固定槽位池设计，避免动态分配导致的复杂组合逻辑：

\begin{itemize}
    \item \textbf{敌机（8 槽，en\_active[7:0]）}：
    轮转生成（en\_spawn\_idx 0$\rightarrow$7$\rightarrow$0），LFSR 随机 X 坐标（10 位），
    生成间隔由 \texttt{ENEMY\_SPAWN\_BASE（120 帧 = 2 秒）}控制。
    每架敌机有独立的移动 AI 和射击模式（共 8 种行为模式），
    子弹写入其专属子槽（敌机 N 对应 eb 槽位 N×8 + 0$\sim$2）。
    敌机下移至屏幕底部（Y $>$ 470）时自动标注 inactive
    \item \textbf{玩家子弹（16 槽，pb\_active[15:0]）}：
    每次射击以 J 键（key\_data[4]）触发。双发模式下分配槽 0 和 1（左右对称，偏移 $\pm$5px），
    武器等级 $\ge$ 2 时增加槽 2（中央第三发）。槽 3-7 实现独立更新。
    子弹以 8 像素/帧速度向上移动，出屏（Y $\le$ 10）标记 inactive
    \item \textbf{敌方子弹（64 槽，eb\_active[63:0]）}：
    每敌机专属 8 槽连续空间。槽 0-7 归敌机 0，槽 8-15 归敌机 1，以此类推。
    子弹以 2 像素/帧速度向下移动，出屏（Y $>$ 470）标记 inactive
\end{itemize}

\paragraph{碰撞检测}

采用逐对 AABB（轴对齐包围盒）硬编码检测，无嵌套 \texttt{for} 循环，确保综合结果可预测：

\begin{itemize}
    \item 玩家子弹 vs 敌机（pb0/1/2 vs en0/1/2/3/4/5/6/7，共约 20 对）：命中后敌机被击毁（HP $\le$ 1 即死），子弹消失，
    总分 +5，击杀数 +1
    \item 敌弹 vs 玩家（eb0/1/2/8/9/10/16/17/18...，共约 19 对，需 pl\_inv==0）：
    命中后 HP 减 1。若 HP $\le$ 1 则进入 GAMEOVER；否则回退至屏幕中央 (320,400)，
    获得 180 帧无敌保护（约 3 秒）+ 清空所有敌弹
    \item 玩家 vs 敌机（en0/1/2/3/4）：命中后 HP 减 1，逻辑同敌弹伤害处理
\end{itemize}

\paragraph{计分与武器升级}

\begin{itemize}
    \item 击毁敌机：+5 分/架（与敌机 HP 无关）
    \item 生存奖励：每 60 帧（1 秒）surv\_cnt 循环 +1 分
    \item 积分模式目标：500 分通关
    \item 武器升级：击杀数达到阈值（10/20/40...，每次翻倍）自动升级（最大 3 级），
    提升子弹数量和缩短射击冷却
\end{itemize}

\paragraph{玩家属性}

\begin{itemize}
    \item 初始位置：(320, 400)（屏幕中央偏下）
    \item 移动速度：2 像素/帧，边界限制 X $\in$ [20, 620]、Y $\in$ [20, 455]
    \item 无敌时间：受伤后 180 帧（约 3 秒），期间 \texttt{pl\_flash} 按 \texttt{pl\_inv[2:1] != 0} 闪烁
    \item 开火冷却：基础 8 帧，随武器等级缩短为 \texttt{FIRE\_RATE - weapon}
    \item 武器升级：击杀 $\ge$ next\_kills\_upg（初始 10，每次翻倍）时触发
\end{itemize}

\paragraph{LFSR 伪随机数生成器}

LFSR 实现在 \texttt{game\_logic.v} 内部（未使用独立模块），
16 位 Galois 型，多项式 $x^{16}+x^{14}+x^{13}+x^{11}+1$（反馈系数 16'hB400）。
每个 game tick 更新一次，种子值 16'hACE1：

\begin{lstlisting}
if (tick)
    lfsr <= {lfsr[14:0], 1'b0} ^ ({16{lfsr[15]}} & 16'hB400);
\end{lstlisting}

LFSR 输出用于：敌机 X 坐标随机化、敌机射击间隔随机扰动、敌机移动方向选择、
障碍物大小/形状/HP 随机化。

\subsubsection{字库实现}

本工程在 \texttt{menu\_render.v} 和 \texttt{hud\_render.v} 中使用 Xilinx Distributed Memory Generator
生成的 ROM IP（\texttt{ROM\_f}）实现 8$\times$8 字库，通过 \texttt{i.coe} 初始化文件（同目录下）
提供位图数据。ROM 地址格式为 \texttt{\{char\_code[5:0], row[2:0]\}}（9 位，512 深度），
8 位数据输出对应一行像素。

该方案替代了旧版的 \texttt{case} 组合逻辑查找表方案，在面积和时序上均有优势。
但 \texttt{hud\_render.v} 中仍保留了一组 \texttt{case} 查找表用于特定字符的 inline 字库。

支持的字符包括 A$\sim$Z（编码 0$\sim$25）、0$\sim$9（编码 26$\sim$35）、
空格（36）、冒号（37）、竖线（39）、短横（40）、大于号（41）、句点（42）。

\subsubsection{渲染器模块族}

所有渲染器（player\_render, enemy\_render, bullet\_render,
hud\_render, menu\_render）均为纯组合逻辑（无 \texttt{posedge}），
基于扫描坐标 (sx, sy) 实时判断每个像素是否属于对应实体。
渲染器无时钟输入端，输出 hit（1 位）和 12 位 BGR 颜色信号（12 位）。

各渲染器的主要几何特征：
\begin{itemize}
    \item \textbf{玩家（player\_render）}：向上等腰三角形飞机，顶点 (px, py-12)、底角 (px$\pm$16, py+12)。
    细节分层：引擎发光（橙色）、驾驶舱（深蓝）、机翼（白色）、主体（青色）。
    支持无敌闪烁（flash=0 时不可见）
    \item \textbf{敌机（enemy\_render）}：椭圆 UFO，简化近似公式 $dx^2 + 2\cdot dy^2 \le 200$（a≈14, b≈10）。
    细节：舷窗（小方形窗口）、穹顶（椭圆弧线）、边缘轮廓（2 像素宽环）。
    HP 分级着色（3→品红、2→紫色、1→暗红），击中闪烁为白色
    \item \textbf{子弹（bullet\_render）}：通过 \texttt{is\_player} 参数区分。
    玩家弹：曼哈顿距离 $\le$ 2 的菱形（青色），敌弹：3$\times$3 方块（红色）
    \item \textbf{菜单（menu\_render）}：双层装饰边框（彩虹色交替 + 浅灰细线）、
    标题文字（AEROPLANE / DANMAKU SHOOTING）、START 提示框、难度和模式显示、操作说明
    \item \textbf{HUD（hud\_render）}：右下角 SCORE 显示（14 字符）、左上角 HP 显示（4 字符）、
    暂停 $\|$$\|$ 符号（黄色）、居中大号 GAME OVER（红色）/ YOU WIN（绿色）
\end{itemize}


% ==================== 四、程序流程 ====================
\section{程序流程}

\subsection{游戏主循环}

系统在 100MHz 时钟域运行。\texttt{vga\_top} 通过三级同步器检测 vs 场同步的下降沿，
在 V-blank 期间产生一个 100MHz 时钟周期的 \texttt{frame\_tick} 脉冲（约 60Hz）。
\texttt{game\_logic} 在每个 frame\_tick 上升沿执行一帧游戏逻辑更新。
这一机制确保了 VGA 画面渲染与游戏逻辑的严格同步，避免画面撕裂。

每帧（frame\_tick 脉冲）执行以下步骤：

\begin{enumerate}[label=(\arabic*)]
    \item 更新 LFSR 伪随机数寄存器
    \item 检测 P/Tab 键边沿触发，更新 pause\_key / btn\_cyc 锁存器
    \item 处理 pause\_cooldown 递减
    \item 根据当前 state 执行对应逻辑：
    \begin{description}
        \item[MENU 状态] 初始化所有实体寄存器→检测按键 1-4 设定难度→检测 M 键设定模式→
        等待 start\_key 进入 PLAYING（写入难度、模式、初始 HP，设置 GRACE\_PERIOD 开局保护）
        \item[PLAYING 状态]\mbox{}
        \begin{itemize}
            \item game\_time 递增，state\_timer 递减
            \item surv\_cnt 循环累计（60 帧 +1 生存分）
            \item 玩家移动（受边界 20$\sim$620 / 20$\sim$455 限制，速度 2px/帧）
            \item 无敌计时（pl\_inv 递减、pl\_flash 闪烁控制）
            \item 玩家射击（空闲槽分配、冷却检查、武器等级决定子弹数）
            \item 移动玩家子弹（向上 8px/帧，出屏标记 inactive）
            \item 敌机生成（轮转空闲槽，LFSR 随机 X，生成间隔 120 帧）
            \item 敌机下漂 + 左右微移 + 敌机专属槽射击
            \item 移动敌方子弹（向下 2px/帧，出屏标记 inactive）
            \item 碰撞检测（逐对 AABB）→ 扣血/击毁/得分/无敌保护
            \item 武器升级检查（total\_kills $\ge$ next\_kills\_upg，自动升级无需按键）
            \item 检查 GAMEOVER（HP=0）或 WIN（score $\ge$ 500 且积分模式）
        \end{itemize}
        \item[PAUSED 状态] 不更新时间，等待 pause\_key 或 fire\_key 恢复（带冷却）
        \item[GAMEOVER/WIN 状态] state\_timer 倒计时（300 帧 → 5 秒）→ 返回 MENU，清除所有实体
    \end{description}
    \item 更新 7 段数码管显示数据（根据 seg\_cycle 选择击杀/生命/积分）
    \item 更新 LED 溢出显示（积分 $\ge$ 10000 或击杀 $\ge$ 100 时点亮高位 LED）
\end{enumerate}

\begin{figure}[H]
\centering
\resizebox{\textwidth}{!}{%
\begin{tikzpicture}[
    node distance = 0.7cm,
    % 竖排文字：窄 text width 令中文逐字换行，纵向排列
    proc/.style = {
        rectangle, draw = black, thick,
        text width = 1.0cm,
        align = center, font = \footnotesize, rounded corners = 2pt,
        fill = blue!8, inner sep = 3pt
    },
    decision/.style = {
        diamond, draw = black, thick,
        aspect = 1.6,
        align = center, font = \footnotesize, fill = yellow!15,
        inner sep = 1pt
    },
    term/.style = {
        rectangle, draw = black, thick,
        text width = 1.1cm,
        align = center, font = \footnotesize, rounded corners = 10pt,
        fill = gray!12, inner sep = 3pt
    },
    stproc/.style = {
        rectangle, draw = black, thick,
        text width = 1.1cm,
        align = center, font = \footnotesize, rounded corners = 2pt,
        fill = red!10, inner sep = 3pt
    },
    arr/.style = {->, >=Stealth, thick},
    larr/.style = {->, >=Stealth, thick, draw=blue!60},
]

% ===== 主流程（中央列 x=0）=====
\node[term] (START) at (0, 0) {frame\\tick\\脉冲};
\node[proc] (LFSR)   at (0, -2.2) {更\\新\\LFSR};
\node[proc] (INPUT)  at (0, -4.4) {P/Tab\\键\\边沿\\锁存};
\node[proc] (DEC)    at (0, -6.6) {cool-\\down\\timer\\递减};
\node[decision] (CHKSTATE) at (0, -8.9) {state?};

\draw[arr] (START) -- (LFSR);
\draw[arr] (LFSR) -- (INPUT);
\draw[arr] (INPUT) -- (DEC);
\draw[arr] (DEC) -- (CHKSTATE);

% ===== 分支总线（从 CHKSTATE 向下引出，避免箭头交叠）=====
\draw[thick] (CHKSTATE.south) -- ++(0,-0.9) coordinate (BUSA);
\draw[thick] (BUSA) -- ++(0,-1.6) coordinate (BUSB);

% ===== MENU（最左 x=-12）=====
\node[stproc] (MENU_INIT) at (-12, -12.5)
    {实\\体\\初\\始\\化};
\node[proc] (MENU_DIFF) at (-12, -15.2)
    {1-4\\选\\难度\\M\\切换\\模式};
\node[decision] (MENU_START) at (-12, -18.0)
    {J\\键?};
\node[stproc] (MENU_WAIT) at (-9.2, -18.0) {等\\待};

\draw[larr] (CHKSTATE.west) -| node[pos=0.25, fill=white, font=\small] {MENU} (MENU_INIT.north);
\draw[arr] (MENU_INIT) -- (MENU_DIFF);
\draw[arr] (MENU_DIFF) -- (MENU_START);
\draw[arr] (MENU_START) -- node[above, font=\footnotesize] {否} (MENU_WAIT);
\draw[arr, rounded corners=6pt] (MENU_WAIT.north) -- ++(0,1.6) -| (MENU_INIT.north east);
\draw[arr, blue!60, thick] (MENU_START.south) -- node[right, font=\footnotesize] {是$\to$PLAY} ++(0,-1.8);

% ===== PLAYING（中偏左 x=-2）=====
\node[stproc] (PLAY_MOVE) at (-2, -12.5)
    {玩\\家\\移动\\WSAD\\无敌\\计时};
\node[proc] (PLAY_SHOOT) at (-2, -15.4)
    {玩家\\射击\\J\\冷却\\4-级};
\node[proc] (PLAY_PBMOVE) at (-2, -18.0)
    {玩家\\子弹\\上移\\8px};
\node[proc] (PLAY_ENEMY) at (-2, -20.6)
    {敌机\\生成\\移动\\射击\\8槽};
\node[proc] (PLAY_EBMOVE) at (-2, -23.4)
    {敌方\\子弹\\下移\\2px};
\node[stproc] (PLAY_COLLISION) at (-2, -26.0)
    {\textbf{碰}\\\textbf{撞}\\检测\\逐对\\AABB};
\node[proc] (PLAY_UPG) at (-2, -28.8)
    {计分\\武器\\升级\\kills\\阈值};
\node[decision] (PLAY_CHKEND) at (-2, -31.6)
    {HP=0?\\score\\$\ge$500?};

\draw[larr] (BUSA) -| node[pos=0.12, fill=white, font=\small] {PLAYING} (PLAY_MOVE.north);
\draw[arr] (PLAY_MOVE) -- (PLAY_SHOOT);
\draw[arr] (PLAY_SHOOT) -- (PLAY_PBMOVE);
\draw[arr] (PLAY_PBMOVE) -- (PLAY_ENEMY);
\draw[arr] (PLAY_ENEMY) -- (PLAY_EBMOVE);
\draw[arr] (PLAY_EBMOVE) -- (PLAY_COLLISION);
\draw[arr] (PLAY_COLLISION) -- (PLAY_UPG);
\draw[arr] (PLAY_UPG) -- (PLAY_CHKEND);

\draw[arr, blue!60, thick] (PLAY_CHKEND.west) -- ++(-1.8,0)
    node[font=\footnotesize, above, text=red!70!black] {\scriptsize HP=0$\to$GO};
\draw[arr, blue!60, thick] (PLAY_CHKEND.east) -- ++(1.8,0)
    node[font=\footnotesize, above, text=green!50!black] {\scriptsize score$\ge$500$\to$WIN};

% ===== PAUSED（右侧 x=6）=====
\node[stproc] (PAUSED_WAIT) at (6, -12.5)
    {冻结\\不更\\新\\实体};
\node[decision] (PAUSED_CHK) at (6, -15.4)
    {P/J键?\\cool-\\down=0?};
\node[stproc] (PAUSED_RES) at (6, -18.2)
    {恢复\\PLAY\\250帧\\冷却};

\draw[larr] (CHKSTATE.east) -| node[pos=0.25, fill=white, font=\small] {PAUSED} (PAUSED_WAIT.north);
\draw[arr] (PAUSED_WAIT) -- (PAUSED_CHK);
\draw[arr] (PAUSED_CHK) -- node[right, font=\footnotesize] {是} (PAUSED_RES);

% ===== GAMEOVER（左下 x=-6.5）=====
\node[stproc] (GOVER_SHOW) at (-6.5, -21)
    {GAME\\OVER\\300帧\\显示};
\node[decision] (GOVER_CHK) at (-6.5, -24.0)
    {timer\\=0?};
\node[stproc] (GOVER_BACK) at (-6.5, -26.8)
    {清\\实体\\$\to$\\MENU};

\draw[arr] (BUSB) -| node[pos=0.28, fill=white, font=\small] {GAMEOVER} (GOVER_SHOW.north);
\draw[arr] (GOVER_SHOW) -- (GOVER_CHK);
\draw[arr] (GOVER_CHK) -- node[right, font=\footnotesize] {是} (GOVER_BACK);

% ===== WIN（右下 x=10）=====
\node[stproc] (WIN_SHOW) at (10, -21)
    {YOU\\WIN\\300帧\\显示};
\node[decision] (WIN_CHK) at (10, -24.0)
    {timer\\=0?};
\node[stproc] (WIN_BACK) at (10, -26.8)
    {清\\实体\\$\to$\\MENU};

\draw[arr] (BUSB) -| node[pos=0.28, fill=white, font=\small] {WIN} (WIN_SHOW.north);
\draw[arr] (WIN_SHOW) -- (WIN_CHK);
\draw[arr] (WIN_CHK) -- node[right, font=\footnotesize] {是} (WIN_BACK);

% ===== 循环返回 =====
\node[term] (END) at (-2, -35.5) {更新\\7段\\LED\\等下\\一tick};

\draw[arr, dashed] (MENU_WAIT.south) |- (END.west);
\draw[arr, dashed] (PLAY_CHKEND.south) -- ++(0,-1.5) -|
    node[font=\footnotesize, pos=0.6, above] {否} (END.north);
\draw[arr, dashed] (PAUSED_CHK.east) -- ++(6.0,0) |- (END.east);
\draw[arr, dashed] (GOVER_CHK.west) -- ++(-3.0,0) |- (END.west);
\draw[arr, dashed] (WIN_CHK.east) -- ++(3.0,0) |- (END.east);

\end{tikzpicture}%
}
\caption{游戏主循环整体流程图（60Hz frame\_tick 驱动）}
\label{fig:flow}
\end{figure}

\subsection{VGA 渲染流程}

VGA 渲染流程与游戏逻辑并行、流水线式运行：

\begin{enumerate}[label=(\arabic*)]
    \item vgac 在 25MHz 像素时钟驱动下产生 (row\_addr, col\_addr) 扫描坐标
    \item 所有渲染器同步接收坐标，纯组合逻辑并行计算 hit 和 color
    \item 优先级 MUX 按固定层级选择最高优先级的有效像素颜色
    \item 输出 12 位 d\_in 到 vgac
    \item vgac 按 VGA 时序驱动 R/G/B/HS/VS 引脚输出到显示器
    \item 在 V-blank 期间，frame\_tick 脉冲驱动游戏逻辑更新一帧
\end{enumerate}

\begin{figure}[H]
\centering
\resizebox{\textwidth}{!}{%
\begin{tikzpicture}[
    node distance = 0.5cm and 0.6cm,
    proc/.style = {
        rectangle, draw = black, thick,
        minimum width = 2.2cm, minimum height = 0.55cm,
        align = center, font = \footnotesize, rounded corners = 2pt,
        fill = green!10
    },
    hproc/.style = {
        rectangle, draw = black, thick,
        minimum width = 2.6cm, minimum height = 0.6cm,
        align = center, font = \footnotesize, rounded corners = 2pt,
        fill = blue!10
    },
    vgac_box/.style = {
        rectangle, draw = black, very thick,
        minimum width = 2.8cm, minimum height = 0.65cm,
        align = center, font = \footnotesize\bfseries, rounded corners = 3pt,
        fill = green!20
    },
    mux_box/.style = {
        rectangle, draw = black, very thick,
        minimum width = 4.2cm, minimum height = 1.4cm,
        align = center, font = \footnotesize, rounded corners = 3pt,
        fill = orange!15
    },
    sync_box/.style = {
        rectangle, draw = black, thick,
        minimum width = 2.8cm, minimum height = 0.6cm,
        align = center, font = \footnotesize, rounded corners = 2pt,
        fill = red!12
    },
    io_box/.style = {
        rectangle, draw = black, thick,
        minimum width = 2.0cm, minimum height = 0.55cm,
        align = center, font = \footnotesize, rounded corners = 2pt,
        fill = yellow!15
    },
    render_box/.style = {
        rectangle, draw = black, thin,
        minimum width = 2.4cm, minimum height = 0.45cm,
        align = center, font = \scriptsize, rounded corners = 2pt,
        fill = green!5
    },
    arr/.style = {->, >=Stealth, thick},
    arr_bus/.style = {->, >=Stealth, thick, draw=green!50!black},
    arr_sync/.style = {->, >=Stealth, thick, draw=red!60},
]

% ---- 上行：VGA 像素渲染管线 ----
\node[io_box] (CLK) {100MHz\\系统时钟};
\node[proc, right=0.5cm of CLK] (DIV) {${\div}4$\\25MHz vga\_clk};
\node[vgac_box, right=0.5cm of DIV] (VGAC) {\textbf{vgac}\\VGA 控制器};

\node[proc, above=0.5cm of VGAC, fill=yellow!12] (SCAN) {row\_addr[8:0]\\col\_addr[9:0]};
\node[proc, below=0.5cm of VGAC, fill=yellow!12] (RDN) {rdn\\消隐控制};
\node[io_box, right=2.0cm of VGAC] (VGAOUT) {RGB[3:0]×3\\HS / VS\\VGA 显示器};

% 渲染器阵列
\node[render_box, right=0.6cm of SCAN, yshift=0.8cm] (R_PL)  {player\_render};
\node[render_box, below=0.12cm of R_PL] (R_EN)  {enemy\_render ×8};
\node[render_box, below=0.12cm of R_EN]  (R_PB)  {bullet\_render ×80\\(16 PB + 64 EB)};
\node[render_box, below=0.12cm of R_PB]  (R_HUD) {hud\_render};
\node[render_box, below=0.12cm of R_HUD] (R_MENU) {menu\_render};

\draw[arr_bus] (VGAC) -- (SCAN);
\draw[arr_bus] (VGAC) -- (RDN);
\draw[arr_bus, shorten >=1pt] (SCAN.east) -- ++(0.25,0) |- (R_PL.west);
\draw[arr_bus, shorten >=1pt] (SCAN.east) -- ++(0.25,0) |- (R_EN.west);
\draw[arr_bus, shorten >=1pt] (SCAN.east) -- ++(0.25,0) |- (R_PB.west);
\draw[arr_bus, shorten >=1pt] (SCAN.east) -- ++(0.25,0) |- (R_HUD.west);
\draw[arr_bus, shorten >=1pt] (SCAN.east) -- ++(0.25,0) |- (R_MENU.west);

% 优先级 MUX
\node[mux_box, right=0.8cm of R_PB, yshift=-0.3cm] (MUX)
    {\textbf{优先级 MUX}\\Menu(MENU) > HUD\\> Player > PBullet\\> EBullet > Enemy > BG};

\draw[arr] (R_PL.east)  -- ++(0.25,0) |- node[font=\tiny, pos=0.55, fill=white] {hit/col} (MUX.north);
\draw[arr] (R_EN.east)  -- (R_EN.east -| MUX.west);
\draw[arr] (R_PB.east)  -- (R_PB.east -| MUX.west);
\draw[arr] (R_HUD.east) -- (R_HUD.east -| MUX.west);
\draw[arr] (R_MENU.east) -- (R_MENU.east -| MUX.west);

\draw[arr_bus] (MUX.east) -- ++(0.35,0) |- node[pos=0.8, font=\tiny, fill=white] {vga\_data[11:0]} (VGAC.east);
\draw[arr] (VGAC) -- (VGAOUT);
\draw[arr] (CLK) -- (DIV);
\draw[arr] (DIV) -- (VGAC);

% ---- 下行：帧同步通路 ----
\node[sync_box, below=1.2cm of VGAC, yshift=-1.2cm] (VS_SYNC)
    {三级同步器\\vs\_sync[2:0]};
\node[sync_box, right=0.8cm of VS_SYNC] (FTICK)
    {frame\_tick 生成\\(vs 下降沿)\\约 60Hz 脉冲};
\node[hproc, right=0.8cm of FTICK] (LOGIC)
    {\textbf{game\_logic}\\游戏状态更新\\(100MHz 域)};

\draw[arr_sync] (VGAC.south) -- ++(0,-0.4) -| node[pos=0.38, font=\tiny, fill=white] {vs} (VS_SYNC.north);
\draw[arr_sync] (VS_SYNC) -- (FTICK);
\draw[arr_sync] (FTICK) -- (LOGIC);

\draw[arr, bend left=30, draw=red!70!black, thick] (LOGIC.north) --
    node[midway, left, font=\footnotesize, text=red!70!black] {实体状态总线} (R_MENU.south);

% ---- 装饰注释 ----
\node[font=\small\bfseries, above=0.2cm of CLK, text=green!50!black]
    {VGA 像素渲染管线（25MHz vga\_clk 域）};
\node[font=\small\bfseries, below=0.2cm of VS_SYNC, text=red!60]
    {帧同步通路（100MHz 主时钟域）};

\end{tikzpicture}%
}
\caption{VGA 渲染管线流程图}
\label{fig:vga_flow}
\end{figure}


% ==================== 五、调试过程分析 ====================
\section{调试过程分析}

\subsection{语法与综合调试}

在 Vivado 中进行语法检查和综合（Synthesis），通过 Messages 窗口检查 Error 和 Warning，
对变量位宽不匹配、未连接端口、多次驱动等问题逐一修正。


\subsection{主要调试问题与解决方案}

\subsubsection{问题 1：玩家飞机位置与移动异常}

\textbf{现象}：玩家飞机初始位置不在屏幕预期位置（中央偏下），移动范围异常。

\textbf{原因}：玩家初始坐标设置错误，边界检查条件不完整。

\textbf{解决方案}：将玩家初始位置设为 (320, 400)（对应 640$\times$480 屏幕中央偏下），
移动边界限制为 X $\in$ [20, 620]、Y $\in$ [20, 455]，考虑飞机半宽（16px）和半高（12px）。

\subsubsection{问题 2：敌机生成和显示不全}

\textbf{现象}：仅有少数敌机槽（索引 0-3）被更新，其余槽位无响应。

\textbf{原因}：敌机生成、移动和射击逻辑只针对前 4 个槽位硬编码，未覆盖全部 8 个槽位。

\textbf{解决方案}：将敌机的生成（en\_spawn\_idx 轮转覆盖 0$\sim$7）、移动和射击逻辑
扩展到全部 8 个敌机槽。前 4 槽保持完整 AI 逻辑，后 4 槽提供基础的生成和下移功能。

\subsubsection{问题 3：子弹发射数量少}

\textbf{现象}：玩家每次只能发射 1$\sim$2 颗子弹。

\textbf{原因}：子弹分配仅检查固定槽位（0 和 1），未充分利用 16 个槽位，且未实现武器升级的第三发。

\textbf{解决方案}：改为优先轮询空闲槽位分配。基础双发分配到槽 0 和 1（左右对称），
武器等级 $\ge$ 2 时增加第三发到槽 2（中央发射）。槽 3-15 作为预留扩展通道。

\subsubsection{问题 4：文字显示异常}

\textbf{现象}：HUD 和菜单中部分字符（冒号、竖线等）显示为空白或乱码。

\textbf{原因}：字符编码寄存器位宽为 [4:0]（5 位），数值 $\ge$ 32 的编码（如空格 36、冒号 37、竖线 39）被截断。

\textbf{解决方案}：将字符编码变量位宽从 [4:0] 扩展为 [5:0]，相关常量从 5'dN 改为 6'dN。
同时将字库从 inline \texttt{case} 查找表升级为 Xilinx Distributed Memory ROM IP，
以 \texttt{i.coe} 初始化文件提供位图数据，提升综合效率。

\subsubsection{问题 5：椭球体渲染鬼影点}

\textbf{现象}：敌机 UFO 的椭圆渲染在距敌机约 256 像素处产生周期性鬼影点（伪像）。

\textbf{原因}：椭圆距离平方计算 (\texttt{dx*dx + 2*dy*dy}) 的值在远离实体时超过位宽上限导致乘法溢出。
具体为 \texttt{adx2} 和 \texttt{ady2} 使用 16 位宽，而 $\texttt{adx}^2$ 理论最大值为 $640^2 = 409600$（需 19 位）。

\textbf{解决方案}：将 \texttt{adx2} 扩展为 22 位宽、\texttt{ady2} 为 22 位、\texttt{dist2} 为 23 位，
阈值比较字面量统一为 23'd 宽度。



% ==================== 六、核心模块模拟仿真 ====================
\section{核心模块模拟仿真结果}

\subsection{键盘驱动模块（Keyboard）仿真}

对 \texttt{Keyboard} 模块进行仿真，验证 PS/2 Set 2 协议的 make/break 解码逻辑：
模块接收 PS/2 时钟和数据信号，解码输出 8 位按键状态 bitmap 和 4 位难度 one-hot 信号。
键盘按下时产生 make 码（含 E0 扩展前缀处理），松开时产生 break 码（F0 前缀 + make 码）。

\begin{figure}[H]
\centering
% TODO: 插入 Keyboard 仿真波形截图
% 建议在 Vivado 中运行 Behavioral Simulation → 添加 Keyboard 为仿真顶层 → 运行仿真 →
% 截图展示：ps2_clk、ps2_data 输入波形和 key_data[7:0]、key_diff[3:0] 输出波形
% 应能观察到按下 W 键（make 码 1D）后 key_data[0] 变为 1
\fbox{\parbox{0.85\textwidth}{\centering\vspace{3cm}【图：Keyboard 仿真波形 - 展示 PS/2 解码效果】\vspace{3cm}}}
\caption{Keyboard 模块仿真波形}
\label{fig:sim_keyboard}
\end{figure}

\subsection{游戏逻辑模块（game\_logic）仿真}

对 \texttt{game\_logic} 模块进行仿真，验证 FSM 状态转换、实体更新、碰撞检测、计分系统等。

主要验证项：
\begin{itemize}
    \item MENU $\rightarrow$ PLAYING 状态转换（start\_key 触发）
    \item PLAYING $\rightarrow$ PAUSED $\rightarrow$ PLAYING 暂停恢复循环
    \item 玩家移动边界限制（X 20$\sim$620, Y 20$\sim$455）
    \item 玩家子弹发射与槽位分配逻辑
    \item 敌机轮转生成、移动与射击
    \item AABB 碰撞检测判定
    \item HP 扣减与无敌保护（pl\_inv 倒计时 + pl\_flash 闪烁）
    \item GAMEOVER 触发条件（HP = 0）
    \item WIN 触发条件（积分模式 score $\ge$ 500）
    \item 武器升级检查（total\_kills $\ge$ next\_kills\_upg）
\end{itemize}

\begin{figure}[H]
\centering
% TODO: 插入 game_logic 仿真波形截图
% 建议展示以下信号：
%   - game_state[2:0]：状态机当前状态
%   - pl_hp[2:0]：玩家生命值变化
%   - total_score[19:0]：积分累计
%   - en_active[7:0]：敌机活跃状态
%   - pb_active[15:0]：玩家子弹活跃状态
%   重点截取一段包含状态转换（MENU→PLAYING→PAUSED→PLAYING→GAMEOVER）的波形
\fbox{\parbox{0.85\textwidth}{\centering\vspace{3cm}【图：game\_logic 仿真波形 - 展示状态转换和关键信号】\vspace{3cm}}}
\caption{game\_logic 模块仿真波形}
\label{fig:sim_logic}
\end{figure}

\subsection{VGA 渲染管线仿真}

上板实测为主，仿真时重点检查 \texttt{vga\_top} 的 RGB 数据时序和同步信号。

\begin{figure}[H]
\centering
% TODO: 插入 vga_top 仿真波形截图
% 建议展示以下信号：
%   - row_addr[8:0], col_addr[9:0]：扫描坐标递进（应看到 0→479 / 0→639 的完整扫描）
%   - rdn：有效区为低 / 消隐区为高
%   - hs, vs：行/场同步脉冲
%   - r[3:0], g[3:0], b[3:0]：像素颜色数据
% 也可附上一段 hit_player / hit_hud 信号展示渲染器的像素命中情况
\fbox{\parbox{0.85\textwidth}{\centering\vspace{3cm}【图：vga\_top 仿真波形 - 展示 VGA 时序和像素输出】\vspace{3cm}}}
\caption{vga\_top VGA 渲染管线仿真波形}
\label{fig:sim_vga}
\end{figure}


% ==================== 七、操作说明 ====================
\section{操作说明}

\subsection{游戏启动}
\begin{enumerate}[label=(\arabic*)]
    \item 将 .bit 文件通过 Vivado Hardware Manager 下载到 K7 FPGA 开发板
    \item 连接 VGA 显示器到开发板 VGA 接口
    \item 按下全局复位键（rstn），系统进入主菜单界面
\end{enumerate}

\subsection{菜单操作}
\begin{itemize}
    \item 按键 1$\sim$4：选择难度（EASY / NORMAL / HARD / HELL）
    \item 按键 M：切换模式（积分 / 无尽）
    \item 按下 J / Space / Enter：开始游戏（前 5 秒为开局保护期，不刷新敌机和障碍物）
\end{itemize}

\subsection{游戏操作}
\begin{itemize}
    \item W/S/A/D：上/下/左/右移动玩家飞机
    \item J：发射子弹（按住连发，受冷却限制）
    \item P：暂停/继续（单次触发，约 4 秒冷却）
    \item Tab：切换 7 段数码管显示内容（击毁数→生命值→积分→循环）
\end{itemize}

\subsection{游戏机制}
\begin{itemize}
    \item 默认 3 条生命（HELL 难度为 2）
    \item 击毁敌机 +5 分，每秒生存奖励 +1 分
    \item 积分模式达到 500 分显示 YOU WIN，5 秒后返回菜单
    \item 无尽模式无通关条件，尽量延长存活时间
    \item 生命归零显示 GAME OVER，5 秒后返回菜单
    \item 受伤后 180 帧（约 3 秒）无敌保护，飞机闪烁提示
    \item 击杀数达 10/20/40 依次升级武器（单发→双发→三发）
\end{itemize}


% ==================== 八、验证过程演示 ====================
\section{验证过程演示}

\subsection{上板验证流程}

\begin{enumerate}[label=(\arabic*)]
    \item Vivado 完成综合（Synthesis）→ 实现（Implementation）→ 生成 .bit 文件（Generate Bitstream）
    \item 通过 Hardware Manager → Program Device 下载 .bit 到 K7 开发板
    \item 连接 VGA 显示器，确认显示正常
    \item 按下全局复位键（rstn），检查主菜单画面
    \item 切换难度（按键 1-4）和模式（M 键），验证菜单文字实时变化
    \item 按下开始键（J/Space/Enter）进入游戏，验证：
    \begin{itemize}
        \item 前 5 秒不刷新敌机和障碍物
        \item 玩家飞机移动范围和速度（WASD）
        \item 子弹发射（双发 + 武器升级后的三发）
        \item 敌机生成、移动和射击
        \item 障碍物生成和漂移
        \item 碰撞检测：玩家子弹击中敌机（得分+5）、击中障碍物（消失）
        \item 玩家受伤：HP-1、无敌闪烁、弹幕清空
        \item HP 归零 → GAMEOVER 画面
        \item 积分 $\ge$ 500（积分模式）→ YOU WIN 画面
    \end{itemize}
    \item 验证暂停/恢复（P 键）
    \item 验证 7 段数码管显示切换（Tab 键）
    \item 验证武器升级（K 键，达到击杀阈值后）
\end{enumerate}

\subsection{上板测试结果}

\begin{figure}[H]
\centering
% TODO: 在此插入上板验证照片（共 6-7 张）
% 建议拍摄以下画面（拍摄时确保画面清晰、反光最小）：
%   1. 主菜单界面（展示标题 AEROPLANE DANMAKU SHOOTING、START 框、难度和模式文字）
%   2. 游戏进行中画面（展示玩家飞机 + 敌机 UFO + 子弹 + 障碍物 + HUD SCORE/HP）
%   3. 暂停界面（展示 || 符号）
%   4. GAMEOVER 界面（展示红色 GAME OVER 大字）
%   5. YOU WIN 界面（展示绿色 YOU WIN 大字）
%   6. 7 段数码管特写（展示 KILLS/LIVES/SCORE 三种模式）
%   7. 整体实验环境照片（K7 开发板 + VGA 显示器 + 电源接线，展示完整实验环境）
%
% 插入方式：每张图使用 subfigure 环境排版，示例：
% \begin{subfigure}{0.48\textwidth}
%     \centering
%     \includegraphics[width=\textwidth]{figs/menu.jpg}
%     \caption{主菜单界面}
% \end{subfigure}
%
\fbox{\parbox{0.85\textwidth}{\centering\vspace{4cm}
【图：上板验证照片合集】

1. 主菜单界面
2. 游戏进行中画面
3. 暂停界面
4. GAMEOVER 界面
5. YOU WIN 界面
6. 7 段数码管特写
7. 整体实验环境
\vspace{4cm}}}
\caption{上板验证结果}
\label{fig:verify}
\end{figure}


% ==================== 九、程序代码（关键部分） ====================
\section{程序代码（关键部分）}

% 完整源码见工程文件，以下仅列出关键模块的核心逻辑。

\subsection{game\_top --- 顶层互联}

\begin{lstlisting}
// 32 位自由运行计数器（100MHz 时钟分频源）
reg [31:0] clkdiv;
always @(posedge clk) clkdiv <= clkdiv + 1'b1;

// PS/2 键盘驱动例化
wire [7:0] key_data;       // 按键状态 bitmap
wire [3:0] key_diff;       // 难度选择 one-hot
wire       key_mode;       // 模式切换
wire       key_cycle;      // 显示切换

Keyboard u_keyboard (
    .clk(clk), .rstn(rstn),
    .ps2_clk(ps2_clk), .ps2_data(ps2_data),
    .key_data(key_data), .key_diff(key_diff),
    .key_mode(key_mode), .key_cycle(key_cycle)
);
\end{lstlisting}

\subsection{game\_logic --- 状态机与输入映射}

\begin{lstlisting}
// Keyboard 模块输出：
// key_data[0]=W(S上), key_data[1]=S(下), key_data[2]=A(左), key_data[3]=D(右)
// key_data[4]=J(开火/开始), key_data[6]=P(暂停)
wire move_up    = key_data[0];   // W
wire move_down  = key_data[1];   // S
wire move_left  = key_data[2];   // A
wire move_right = key_data[3];   // D
wire fire_key   = key_data[4];   // J
wire start_key  = fire_key;
\end{lstlisting}

\subsection{game\_logic --- FSM 状态转换}

\begin{lstlisting}
case (state)
`STATE_MENU: begin
    // 初始化所有实体寄存器
    pl_alive <= 1'b1;  pl_x <= 10'd320;  pl_y <= 9'd400;
    total_score <= 20'd0;  total_kills <= 8'd0;
    en_act <= 8'd0;  pb_act <= 16'd0;  eb_act <= 64'd0;

    // 难度与模式选择
    if (key_diff[0]) sel_diff <= `DIFF_EASY;    // 按键 1
    if (key_diff[1]) sel_diff <= `DIFF_NORMAL;  // 按键 2
    if (key_diff[2]) sel_diff <= `DIFF_HARD;     // 按键 3
    if (key_diff[3]) sel_diff <= `DIFF_HELL;     // 按键 4
    sel_mode <= sel_mode ^ key_mode;              // M 键切换

    if (start_key) begin
        state       <= `STATE_PLAYING;
        state_timer <= `GRACE_PERIOD;  // 300 帧 = 5 秒开局保护
        difficulty  <= sel_diff;
        mode        <= sel_mode;
        pl_hp       <= (sel_diff == `DIFF_EASY)  ? 3'd3 :
                       (sel_diff == `DIFF_HELL)  ? 3'd2 : 3'd3;
    end
end

`STATE_PLAYING: begin
    // 暂停检测（带冷却）
    if (pause_key && pause_cooldown == 8'd0) begin
        state <= `STATE_PAUSED;
        pause_cooldown <= 8'd250;  // 约 4 秒冷却
    end

    // …… 实体更新、碰撞检测、计分 ……

    // 状态转移检测
    if (pl_hp == 0) begin
        state <= `STATE_GAMEOVER;
        state_timer <= `GAMEOVER_DISPLAY;
    end else if (!mode && total_score >= `SCORE_WIN_TARGET) begin
        state <= `STATE_WIN;
        state_timer <= `WIN_DISPLAY;
    end
end

`STATE_PAUSED: begin
    if ((pause_key || fire_key) && pause_cooldown == 8'd0) begin
        state <= `STATE_PLAYING;
        pause_cooldown <= 8'd250;
    end
end

`STATE_GAMEOVER: begin
    if (state_timer > 0) state_timer <= state_timer - 11'd1;
    else state <= `STATE_MENU;  // 5 秒后返回菜单
end

`STATE_WIN: begin
    if (state_timer > 0) state_timer <= state_timer - 11'd1;
    else state <= `STATE_MENU;  // 5 秒后返回菜单
end
endcase
\end{lstlisting}

\subsection{LFSR 伪随机数生成}

\begin{lstlisting}
// 16 位 Galois LFSR，多项式 x^16 + x^14 + x^13 + x^11 + 1
// 反馈系数：16'hB400 = 16'b1011_0100_0000_0000
// 种子值：16'hACE1
if (tick)
    lfsr <= {lfsr[14:0], 1'b0} ^ ({16{lfsr[15]}} & 16'hB400);
\end{lstlisting}

\subsection{hud\_render --- 内联字库（部分）}

\begin{lstlisting}
// 字符 ROM 综合为 Distributed Memory ROM IP（ROM_f）
// 地址格式：{char_code[5:0], sy[2:0]} = 9 位，512 深度
// 数据位宽：8 位（一行像素的位图）
ROM_f u_font_rom (
    .a  ({char_code[5:0], sy[2:0]}),
    .spo(char_bitmap)
);
// 字库初始化文件：i.coe（项目根目录）
\end{lstlisting}


% ==================== 十、组内成员分工说明及贡献比例 ====================
\section{组内成员分工说明及贡献比例}

\subsection{成员分工}

% TODO: 根据实际情况填写分工

\begin{table}[H]
\centering
\caption{成员分工表}
\label{tab:division}
\begin{tabular}{lp{4.5cm}p{5cm}}
\toprule
\textbf{姓名} & \textbf{负责内容} & \textbf{具体工作} \\
\midrule
{[姓名1]}（组长） & 系统架构、VGA 渲染管线 &
    整体架构设计、game\_top 顶层互联、VGA 渲染管线（vga\_top + vgac + 优先级 MUX）、
    渲染器族（player/enemy/bullet/menu/hud\_render）、字库 ROM、报告撰写 \\
\midrule
{[姓名2]} & 游戏逻辑、外设驱动 &
    game\_logic 核心 FSM、实体管理与碰撞检测、LFSR 随机数、
    PS/2 键盘驱动（Keyboard）、7 段数码管驱动、上板调试 \\
\midrule
{[姓名3]} & 仿真验证、操作说明 &
    各模块单元仿真、上板验证测试、操作说明书编写、实验报告排版 \\
\bottomrule
\end{tabular}
\end{table}

\subsection{贡献比例}

% TODO: 根据实际情况填写比例
\begin{itemize}
    \item {[姓名1]}（组长）：34\%
    \item {[姓名2]}：33\%
    \item {[姓名3]}：33\%
\end{itemize}


% ==================== 十一、总结与展望 ====================
\section{总结与展望}

\subsection{项目总结}

本工程在 Sword Kintex-7 FPGA 平台上实现了一个完整的弹幕射击游戏，涵盖以下数字逻辑设计要点：

\begin{enumerate}[label=(\arabic*)]
    \item \textbf{有限状态机（FSM）设计}：5 状态 Moore 型状态机（MENU/PLAYING/PAUSED/GAMEOVER/WIN），
    完整的状态转换条件判断、边沿触发锁存消费机制和冷却保护
    \item \textbf{VGA 显示控制}：640$\times$480 @ 60Hz，12 位 BGR 色深，纯组合逻辑渲染器族 + 优先级合成管线，
    多实体并行渲染无帧缓冲设计
    \item \textbf{人机交互接口}：PS/2 键盘输入（Keyboard 模块解码），
    VGA 视频 + 7 段数码管（直连 + 串行 P2S）+ LED 多种输出
    \item \textbf{实体系统管理}：8 敌机 + 16 玩家子弹 + 64 敌方子弹，
    固定槽位池管理 + 专属子槽分配，通过扁平化高位宽总线与 generate 展开在综合层面优化
    \item \textbf{碰撞检测}：逐对 AABB 轴对齐包围盒硬编码检测，无嵌套 for 循环，
    多类型交互（玩家子弹 $\times$ 敌机、敌弹 $\times$ 玩家、玩家 $\times$ 敌机）
    \item \textbf{伪随机数生成}：16 位 Galois LFSR，多项式 $x^{16}+x^{14}+x^{13}+x^{11}+1$，
     驱动敌机坐标、射击间隔、移动方向的随机化
    \item \textbf{PS/2 键盘解码}：Keyboard 模块实现 PS/2 Set 2 协议 make/break 解码、E0 扩展前缀处理、8 位按键 bitmap + 4 位难度 one-hot 输出
    \item \textbf{时钟域设计}：100MHz 系统时钟 → 25MHz VGA 时钟 + 60Hz frame\_tick
    通过三级同步器实现跨时钟域 vs 下降沿检测
    \item \textbf{字库与 ROM}：Xilinx Distributed Memory ROM IP 实现 8$\times$8 字库，
    i.coe 初始化文件提供位图数据，支持 A-Z、0-9 和标点共 40+ 字符
\end{enumerate}

\subsection{改进方向}

\begin{enumerate}[label=(\arabic*)]
    \item \textbf{音效系统}：开发板蜂鸣器（Buzzer）可接入简单 PWM 驱动，
    实现射击音效、爆炸音效、背景音乐（BGM）等音频反馈
    \item \textbf{曲线弹道}：当前敌方子弹均为直线弹道（BTYPE\_STRAIGHT），
    \texttt{game\_defs.vh} 已预留抛物线（BTYPE\_PARABOLA）、顺时针圆弧（BTYPE\_CIRCLE\_CW）、
    逆时针圆弧（BTYPE\_CIRCLE\_CCW）三种曲线类型编码，可在 \texttt{game\_logic} 中实现对应的运动方程更新
    \item \textbf{敌机 AI 多样化}：当前 8 个敌机槽各有不同的移动和射击行为模式，
    可在全部 8 个槽位实现不同 AI 策略（悬停射击、S 形移动、追踪玩家、弹幕散布等）
    \item \textbf{图形质量提升}：当前玩家飞机和敌机 UFO 为几何图形（三角形/椭圆），
    可通过 Block RAM 存储精灵位图（sprite）实现更精细的飞机和 UFO 造型
    \item \textbf{难度动态调整}：\texttt{game\_defs.vh} 已定义难度系数（DMUL\_EASY = 1.0、
    DMUL\_NORMAL = 1.5、DMUL\_HARD = 2.0、DMUL\_HELL = 3.0）和时间增长函数 T(t) 的框架，
    可在 \texttt{game\_logic} 中实现随时间动态调整敌机刷新间隔、HP 和子弹密度
\end{enumerate}


% ==================== 参考资料 ====================
\section{参考资料}

\begin{enumerate}[label={[\arabic*]}]
    \item NJU 数字逻辑设计实验讲义，PS/2 键盘输入实验，
          \url{https://nju-projectn.github.io/dlco-lecture-note/exp/07.html}
    \item YooLc/FPGA-STG：基于 FPGA 的 STG 射击游戏，
          \url{https://github.com/YooLc/FPGA-STG}
    \item Xilinx, ``Vivado Design Suite User Guide: Synthesis,'' UG901
    \item Xilinx, ``7 Series FPGAs Data Sheet: Overview,'' DS180
    \item Xilinx, ``Vivado Design Suite User Guide: I/O and Clock Planning,'' UG899
    \item J. Bhasker 著，徐振林等译，``Verilog HDL 硬件描述语言''，机械工业出版社，2000
    \item 夏宇闻，``Verilog 数字系统设计教程''，北京航空航天大学出版社
\end{enumerate}

\end{document}
